\documentclass[preprint,12pt,authoryear]{elsarticle}
\usepackage{amsmath,amssymb,amsfonts,amsthm}
\usepackage{graphicx}
\usepackage{booktabs}
\usepackage{tabularx}
\usepackage{geometry}
\usepackage{natbib}
\usepackage{hyperref}

\geometry{a4paper, margin=1in}

\newtheorem{theorem}{Theorem}
\newtheorem{proposition}{Proposition}
\newtheorem{corollary}{Corollary}
\newtheorem{lemma}{Lemma}
\theoremstyle{definition}
\newtheorem{definition}{Definition}

\journal{Ecological Economics}

\begin{document}

\begin{frontmatter}

\title{The Physics of Economic Value: Thermodynamic Exergy, Metabolic Labor, and Stock-Flow Consistent Dynamic Modeling of Subsistence Supply Chains}

\author{Jarrod Lyndsay Hamilton}
\address{Alethekanon Research Institute, Australia}

\begin{abstract}
Standard macroeconomic models formalize price as a single-dimensional scalar cleared instantaneously by Walrasian equilibrium ($dt = 0$), assuming frictionless substitution, symmetrical transaction agency, and substrate-blind accounting. These assumptions conceal thermodynamic exergy destruction, human metabolic labor exhaustion, asymmetric survival coercion, and working capital debt compounding.

  

This paper develops a unified mathematical and empirical framework for \textbf{Value Physics}. We formulate the \textbf{Universal Price Equation (UPE)}, decomposing total cost into a Three-Layer substrate stack ($P_t = P_e + P_b + P_m$), where thermodynamic exergy destruction ($P_e$) and biomechanical metabolic effort ($P_b$) form an irreducible physical cost floor ($P_H$). We derive the relativistic Lorentz Coercion Factor ($\gamma(v_{\text{rel}})$), proving from Stone-Geary subsistence utility and Rubinstein bargaining that under existential urgency, demand price elasticity collapses ($\epsilon \to 0$), enabling unearned rentier toll extraction ($R_a \to 1.0$). We construct the 7-Plane Price Tensor ($\mathbf{P}^7$) and the multi-commodity Leontief-Value Tensor $(\mathbf{I} - \mathbf{A})^{-1}$ across five vital human provisioning sectors (Food, Energy, Water, Shelter, Care).

  

Using an empirical \textbf{Stock-Flow Consistent (SFC)} dynamic model of an Australian broadacre agricultural community calibrated against ABARES, ABS, RBA, ACCC, and AEMO datasets across 156 weeks of compound macroeconomic and climatic shocks (RBA cash rate hike +425 bps, diesel fuel increase to \$2.25/L, 40\% drought yield shock), we prove that central bank monetary tightening operates as a cost-push inflationary driver on non-discretionary subsistence staples via the upstream working capital interest channel ($P\_{Q6}$). Finally, we formalize **Theorem 5 (Statutory Debt Bifurcation)** and the **Tri-Fold Cryptographic Ledger**, demonstrating that replacing exponential compound debt ($L\_0 e^{rt}$) with linear principal amortization ($\\frac{C\_{\\text{principal}}}{T\_{\\text{term}}}$) at 0.0\% interest guarantees 100\% solvency across 1,000 Monte Carlo shock iterations while preserving physical carrying capacity ($R\_n$).
\end{abstract}

\begin{keyword}
Econophysics \sep Value Physics \sep Stock-Flow Consistent Modeling \sep Exergoeconomics \sep Cost Channel of Monetary Policy \sep Proof-of-Difficulty \sep Non-Usurious Credit \sep Universal Price Equation
\JEL B52 \sep E12 \sep E43 \sep E52 \sep Q01 \sep Q57
\end{keyword}

\end{frontmatter}

\section{Introduction \& Ontological Critique of Single-Scalar Price Theory}

Neoclassical and New Keynesian economic frameworks formalize price as a one-dimensional equilibrium scalar ($P \in \mathbb{R}^+$) cleared at the intersection of demand and supply curves. While analytically tractable, this compression embeds four fundamental ontological errors:

  

1.  \textbf{Dimensional Compression:} Compressing multi-dimensional biological, physical, logistical, and historical realities into a single nominal fiat scalar masks where real costs are paid—whether by biological human tissue, degraded soil carbon, depleted water tables, or public infrastructure.
2.  \textbf{Space and Time Neutrality:} Assuming instantaneous market clearing ($dt = 0$) at a dimensionless point violates the finite speed of physical causality ($c$) and ignores logistical friction and perishable lifecycle decay.
3.  \textbf{Bilateral Symmetry Fiction:} Treating market exchange as an uncoerced bilateral negotiation between symmetric agents ignores existential urgency differentials, where one agent transacts under immediate biological survival necessity ($v_{\text{rel}} \to 1.0$) while the counterparty possesses infinite temporal liquidity patience.
4.  \textbf{Substrate Blindness:} Treating monetary issuance as an unconstrained symbolic claim disconnected from the physical well-depth ($R_n$) and exergy destruction ($P_e$) required to reproduce the biosphere.

  

Building on the bioeconomic foundations of Georgescu-Roegen (1971), Soddy's (1926) thermodynamic critique of compound interest, exergoeconomics (Tsatsaronis 1993; Valero 2006), and Sraffian standard commodities (Sraffa 1960), this paper formalizes the physics of economic value and integrates it into empirical dynamic macroeconomic models.

\section{Theoretical Architecture of Value Physics}

\subsection{The Universal Price Equation (UPE) \& Three-Layer Cost Decomposition}

Total price $P_t(x, t, c)$ across spacetime and physical containment is decomposed into three distinct layers:  
$$P_t = P_e + P_b + P_m$$  
Where:

  

  - $P_e$ is primary thermodynamic exergy destroyed (measured in $\text{MJ}$ or $\text{kWh}$).
  - $P_b$ is human metabolic and cognitive labor expended (measured in $WEST$ difficulty tokens).
  - $P_m$ is nominal market mediation and financial toll markup.

  

The continuous dynamical micro-foundations of kinetic exergy fluxes ($W_a, A_a$) and structural floor sedimentation ($W_p, A_p$) are formally derived in the companion theoretical treatise on Thermodynamic Work-Antiwork State-Space Dynamics (Hamilton, 2026, ARI-VPE-TH-2026-02), establishing the Autonomous Sclerosis Blowup Theorem and the 45-degree Diamond Phase Space.

The master Universal Price Equation is formulated as:  
$$P(x, t) = m_1 m_2 \gamma(v_{\text{rel}}) \left[ \frac{S(x, t, c) \cdot U(x, t, c)}{R_n(x, t)(1 - R_a(x, t))} \right] + P_e(x, t) + P_b(x, t)$$  
Where $S(x, t, c)$ and $U(x, t, c)$ are spatiotemporal scarcity and urgency fields, $R_n$ is ecological carrying capacity, $R_a \in [0, 1)$ is the rentier enclosure toll fraction, and $\gamma(v_{\text{rel}})$ is the relativistic coercion factor.

\subsection{Relativistic Coercion \& Stone-Geary Elasticity Collapse}

When transacting parties experience asymmetric existential urgency ($U_A \gg U_B$), relative urgency velocity is defined as:  
$$v_{\text{rel}} = \frac{|U_A - U_B|}{\max(U_A, U_B) + 1\infty_U} \in [0, 1), \quad \gamma(v_{\text{rel}}) = \frac{1}{\sqrt{1 - v_{\text{rel}}^2}}$$

  

\begin{theorem}[Theorem 2.1 (Elasticity Collapse under Urgent Coercion):] Grounded in Stone-Geary (Geary 1950; Stone 1954) subsistence utility and Rubinstein (1982) bargaining dynamics, as consumption approaches the subsistence baseline ($v_{\text{rel}} \to 1.0$), demand price elasticity collapses to zero:  
$$\epsilon(v_{\text{rel}}) = \epsilon_0 \sqrt{1 - v_{\text{rel}}^2} \xrightarrow{v_{\text{rel}} \to 1.0} 0$$  
Price sensitivity vanishes completely, enabling monopolistic rentier toll extraction ($R_a \to 1.0$) on captive populations without reducing unit volume (formal mathematical derivation provided in Appendix B).
\end{theorem}

\subsection{The 7-Plane Price Tensor ($\mathbf{P}^7$) \& Value Gradient}

Price is an irreducible 7-dimensional tensor across orthogonal reality planes:  
$$\mathbf{P}^7 = \begin{pmatrix} P_{Q1} & (\text{WHO: Identity / Metaphysical Access}) \\ P_{Q2} & (\text{WHAT: Opportunity / Alternative Foregone}) \\ P_{Q3} & (\text{WHERE: Logistics / Transport Friction}) \\ P_{Q4} & (\text{WHY: Prestige / Symbolic Markup}) \\ P_{Q5} & (\text{HOW: Operational / Physical Transformation}) \\ P_{Q6} & (\text{CAUSE: Historical Sunk Capital \& Debt Service}) \\ P_{Q7} & (\text{EFFECT: Consequential / Ecological Remainder}) \end{pmatrix}$$  
Economic value flows along the negative gradient of systemic price strain on the Riemannian bundle manifold $\mathcal{M}^7$:  
$$\mathbf{V}^7 = -\nabla_{\mathcal{M}^7} \mathbf{P}^7$$

\subsection{Theorem 2.4 (Kantorovich-Samuelson Spatial Value Duality)}

\begin{theorem}[Theorem 2.4 (Kantorovich-Samuelson Spatial Value Duality):] Let the spatial transport cost between geographic nodes $i$ and $j$ along plane $P_{Q3}$ be $c_{ij} = D_{ij} \cdot \kappa_{\text{exergy}}$, where $D_{ij}$ is the network transport distance and $\kappa_{\text{exergy}}$ is the physical exergy destruction rate per unit distance. Under the Monge-Kantorovich minimum exergy transportation problem:

$$\min \sum_{i,j} c_{ij} x_{ij} \quad \text{subject to } \sum_j x_{ij} \le S_i, \quad \sum_i x_{ij} \ge D_j, \quad x_{ij} \ge 0$$

The dual linear programming formulation yields origin supply shadow potentials $u_i^*$ and destination delivery shadow prices $v_j^*$. By the Strong Duality Theorem of Linear Programming and the Karush-Kuhn-Tucker (KKT) complementary slackness conditions, on every active physical trade route ($x_{ij}^* > 0$):

$$v_j^* - u_i^* = c_{ij} = D_{ij} \cdot \kappa_{\text{exergy}}$$

While for all inactive trade routes:

$$v_j^* - u_i^* \le c_{ij}$$

Consequently, the Directional Value Gradient $\mathbf{V}^7 = -\nabla_{\mathcal{M}^7} \mathbf{P}^7$ evaluated along the spatial logistics dimension ($P_{Q3}$) corresponds exactly to the optimal dual shadow price gradient:

$$\nabla_{ij} P_{Q3} = v_j^* - u_i^* = c_{ij}$$

Under non-usurious clearing ($R_a \equiv 0$), spatial arbitrage profits collapse to zero:

$$\Pi_{\text{arbitrage}} = (P_j - P_i) - c_{ij} \equiv 0$$

Regional price differentials equate strictly to the irreversible thermodynamic work of transport ($P_e$).
\end{theorem}

\subsection{The WEST Token Standard \& The Distortion Quotient ($DQ$)}

Instantaneous labor difficulty is formalized as:  
$$D(t) = s(t) \cdot e(t) \cdot [1 + d(t)] \cdot [1 + c(t)]$$  
Where $s(t) = 1 + \ln(1 + \text{Training Years})$ is skill density, $e(t)$ is metabolic exertion in METs, $d(t)$ is physical danger, and $c(t)$ is cognitive focus density. $1\text{ WEST Token} \equiv 1\text{ hour}$ of labor at $s=1.0, e=1.0\text{ MET}, d=0.0, c=0.0$.

  

The Distortion Quotient ($DQ$) measures institutional rentier extraction across supply chains:  
$$DQ = \frac{\% \text{ Share of Nominal Market Price } (P_m)}{\% \text{ Share of Physical Difficulty } (WEST)}$$

  

  - $DQ = 1.00$: True physical difficulty parity.
  - $DQ \ll 1.00$: Physical substrate exploitation.
  - $DQ \gg 1.00$: Unearned rentier toll capture.

  

\subsection{The Invertible Algebraic State Equation \& Diagnostic Measurement Operators (Logarithmic Arithmetic Polarity)}

Standard neoclassical price theory treats price as an exogenous scalar cleared by market equilibria. In Value Physics, the Universal Price Equation (UPE) operates as an invertible algebraic equation of state—analogous to the ideal gas law ($PV = nRT$) in thermodynamics or equations of motion ($F = ma$) in Newtonian mechanics.

Decomposing the kinetic markup term $P_{\text{kinetic}} = P_t - (P_e + P_b)$ into its logarithmic representation reveals an algebraic product group with discrete arithmetic polarities:

$\ln(P_{\text{kinetic}}) = [ +\ln(m_1) + \ln(m_2) + \ln(\gamma) + \ln(S) + \ln(U) ] - [ \ln(R_n) + \ln(1 - R_a) ]$

Where parameters with positive arithmetic polarity ($+1$) act as numerator price multipliers, and parameters with negative arithmetic polarity ($-1$) act as denominator capacity dampeners. Sunk thermodynamic exergy ($P_e$) and metabolic labor ($P_b$) serve as invariant additive translations anchoring the physical floor.

Because the state equation is algebraically invertible across all parameters, the UPE functions as an empirical diagnostic measurement instrument. By observing market transaction prices $P_t$ alongside physical exergy expenditures ($P_e$) and labor ($P_b$), economists can invert the equation to solve for unobservable institutional and biophysical variables:

1.  1. The Monopoly Enclosure Detector ($R_a$): Isolates the percentage of market price siphoned by predatory gatekeeping and institutional tolls:  
    $R_a = 1 - [ (m_1 \cdot m_2 \cdot \gamma(v_{\text{rel}}) \cdot S \cdot U) / (R_n \cdot (P_t - P_e - P_b)) ]$
2.  2. The Substrate Carrying Capacity Metric ($R_n$): Measures the physical resource well-depth and ecological degradation from observed price strain:  
    $R_n = [ (m_1 \cdot m_2 \cdot \gamma(v_{\text{rel}}) \cdot S \cdot U) / ((P_t - P_e - P_b) \cdot (1 - R_a)) ]$
3.  3. The Survival Coercion Meter ($v_{\text{rel}}$): Calculates consumer existential desperation directly from observed price markups above physical cost:  
    $v_{\text{rel}} = \sqrt{ 1 - [ (m_1 \cdot m_2 \cdot S \cdot U) / ((P_t - P_e - P_b) \cdot R_n \cdot (1 - R_a)) ]^2 }$
4.  4. The Metabolic Labor Floor ($P_b$): Strips away speculative rents and scarcity premiums to measure pure physical human exertion:  
    $P_b = P_t - P_e - m_1 \cdot m_2 \cdot \gamma(v_{\text{rel}}) \cdot [ (S \cdot U) / (R_n \cdot (1 - R_a)) ]$

This algebraic invertibility transforms Value Physics from a theoretical pricing framework into an operational diagnostic suite for economic auditors and regulators.

\subsection{The Decoupled Identification Theorem for Institutional Tolls vs. Ecological Depletion (Theorem 2.7)}

While Section 2.6 establishes the analytical inversion for any single unobserved parameter, an identification problem arises when both monopoly enclosure ($R_a(t)$) and ecological carrying capacity ($R_n(t)$) are simultaneously unknown. Because they enter multiplicatively in the denominator of the kinetic term, an increase in consumer price can be caused either by biophysical resource depletion ($R_n \to 0$) or by predatory supermarket markup expansion ($R_a \to 1.0$).

We solve this identification challenge across multi-period panel data ($t = 1, \dots, T$) by coupling the economic price observation equation to the physical transition equation of ecological capital.

Let normalized observable log-strain be defined as:  
$Y(t) = \ln(P_t(t) - P_e(t) - P_b(t)) - \ln(m_1(t) m_2(t) \gamma(v_{\text{rel}}(t)) S(t) U(t)) = -\ln R_n(t) - \ln(1 - R_a(t))$

Unlike financial claims, physical substrate carrying capacity $R_n(t)$ (such as topsoil organic carbon, dam water volume, or processing infrastructure) is governed by physical conservation laws subject to measurable physical investment $I_n(t)$ and depreciation $\delta_n$:  
$R_n(t+1) = (1 - \delta_n) R_n(t) + \alpha_n I_n(t)$

Taking the forward difference of log-strain yields:  
$\Delta Y(t) = Y(t+1) - Y(t) = -\Delta \ln R_n(t) - \Delta \ln(1 - R_a(t))$

\begin{theorem}[Theorem 2.7 (Decoupled Identification of Institutional Rents and Biophysical Depletion):] Because the biophysical trajectory $\Delta \ln R_n(t) = \ln[ (1 - \delta_n) + \alpha_n (I_n(t) / R_n(t)) ]$ is uniquely determined by observable physical capital accounts, the change in institutional monopoly extraction $\Delta \ln(1 - R_a(t))$ is uniquely and orthogonally identified:

$\Delta \ln(1 - R_a(t)) = -\Delta Y(t) - \Delta \ln R_n(t)$

$R_a(t+1) = 1 - (1 - R_a(t)) \cdot \exp\left( -[\Delta Y(t) + \Delta \ln R_n(t)] \right)$
\end{theorem}

\begin{corollary}[Corollary 2.7 (Orthogonal Audit Resolution):]

1.  1\. If a price increase is driven purely by physical resource depletion (e.g. severe drought where $I_n < \delta_n R_n$ and $\Delta \ln R_n < 0$), then $\Delta Y(t) = -\Delta \ln R_n(t)$, and $R_a(t)$ remains constant.
2.  2\. If a price increase is driven by corporate price gouging, then $\Delta Y(t) > 0$ while capital telemetry confirms $\Delta \ln R_n(t) \ge 0$, uniquely isolating the growth of unearned monopoly extraction $\Delta R_a(t) > 0$.

This theorem provides anti-trust regulators and central banks with an empirical mechanism to distinguish true supply shocks from artificial rent extraction.
\end{corollary}

\subsection{The Stochastic Extended Kalman Filter (EKF) State-Space Observer \& Asymptotic Error Bounds (Theorem 2.8)}

While Section 2.7 establishes the analytical decoupling under deterministic conditions, empirical telemetry from agricultural sensor networks (soil moisture arrays, grain yield monitors) and retail checkout scanner data inherently contains stochastic measurement noise and unmodeled disturbances. 

We formulate a non-linear state-space Extended Kalman Filter (EKF) observer to track the latent state vector:  
$$\mathbf{x}_t = \begin{bmatrix} z_t \\\\ \theta_t \end{bmatrix} = \begin{bmatrix} \ln R_{n,t} \\\\ \ln(1 - R_{a,t}) \end{bmatrix}$$

The non-linear physical state transition model is:  
$$z_{t+1} = z_t + \ln\left( (1 - \delta_n) + \alpha_n \frac{I_{n,t}}{\exp(z_t)} \right) + w_{z,t}$$  
$$\theta_{t+1} = \theta_t + w_{\theta,t}$$

Where $\mathbf{w}_t \sim \mathcal{N}(\mathbf{0}, \mathbf{Q})$ is process noise with covariance $\mathbf{Q} = \text{diag}(q_z, q_\theta)$.

The scalar observation equation for normalized log-strain is:  
$$y_t = \ln(P_{\text{kinetic},t}) - \ln(m_1 m_2 \gamma S U) = -z_t - \theta_t + v_t, \quad v_t \sim \mathcal{N}(0, R)$$  
with linear measurement operator $\mathbf{H} = \begin{bmatrix} -1 & -1 \end{bmatrix}$.

The state transition Jacobian $\mathbf{F}_t = \frac{\partial \mathbf{f}}{\partial \mathbf{x}}$ evaluated at the prior estimate is:  
$$F_{11} = 1 - \frac{\alpha_n I_{n,t} \exp(-z_t)}{(1 - \delta_n) + \alpha_n I_{n,t} \exp(-z_t)}, \quad F_{12} = 0, \quad F_{21} = 0, \quad F_{22} = 1$$

The recursive state estimation proceeds via standard predictor-corrector equations:  
1\. State \& Covariance Prediction:  
$$\hat{\mathbf{x}}_{t|t-1} = \mathbf{f}(\hat{\mathbf{x}}_{t-1|t-1}, I_{n,t-1})$$  
$$\mathbf{P}_{t|t-1} = \mathbf{F}_{t-1} \mathbf{P}_{t-1|t-1} \mathbf{F}_{t-1}^T + \mathbf{Q}$$

2\. Innovation \& Measurement Update:  
$$\tilde{y}_t = y_t - \mathbf{H} \hat{\mathbf{x}}_{t|t-1}$$  
$$S_t = \mathbf{H} \mathbf{P}_{t|t-1} \mathbf{H}^T + R$$  
$$\mathbf{K}_t = \mathbf{P}_{t|t-1} \mathbf{H}^T S_t^{-1}$$  
$$\hat{\mathbf{x}}_{t|t} = \hat{\mathbf{x}}_{t|t-1} + \mathbf{K}_t \tilde{y}_t$$  
$$\mathbf{P}_{t|t} = (\mathbf{I} - \mathbf{K}_t \mathbf{H}) \mathbf{P}_{t|t-1}$$

\begin{theorem}[Theorem 2.8 (Stochastic State-Space Observer Convergence \& Asymptotic Covariance Bound):] Because the linearized pair $(\mathbf{F}_t, \mathbf{H})$ is uniformly observable across all non-zero capital investment regimes ($I_{n,t} > 0$), the discrete-time Extended Kalman Filter observer is exponentially stable in the mean square. The posterior error covariance is strictly bounded:  
$$\lim_{t \to \infty} \text{Tr}(\mathbf{P}_{t|t}) \le \frac{R \cdot \text{Tr}(\mathbf{Q})}{R + \text{Tr}(\mathbf{Q})} < \infty$$
\end{theorem}

\textbf{Empirical Calibration \& Noise Resilience:}

\subsection{High-Order Unscented State Reconstruction \& Curvature Invariance (Theorem 2.9)}

While Theorem 2.8 establishes first-order Extended Kalman Filter (EKF) observer convergence, socio-ecological economies frequently undergo acute non-linear bifurcations—such as rapid climatic drought onsets, bank credit freezes, or sudden regulatory structural shifts. Under such high-curvature dynamics, Jacobian linearization $$\mathbf{F}_t = \left. \frac{\partial \mathbf{f}}{\partial \mathbf{x}} \right|_{\hat{\mathbf{x}}}$$ introduces non-negligible second- and higher-order Taylor truncation errors ($$\mathcal{O}(\\|\Delta \mathbf{x}\\|^2)$$) that can induce filter bias or covariance under-estimation.

To guarantee high-order estimation precision, we formulate the Unscented Kalman Filter (UKF) observer using the deterministic sigma-point architecture of Julier \& Uhlmann (1997, 2004). For an $n$-dimensional state vector ($n=2$ for $$\mathbf{x}_t = [\ln R_{n,t}, \ln(1 - R_{a,t})]^T$$), $2n+1 = 5$ deterministic sigma points $$\boldsymbol{\chi}_i$$ are generated symmetrically around the posterior mean:

$$\boldsymbol{\chi}_0 = \hat{\mathbf{x}}_{t|t}$$

$$\boldsymbol{\chi}_i = \hat{\mathbf{x}}_{t|t} + \left(\sqrt{(n + \lambda_{\text{ukf}})\mathbf{P}_{t|t}}\right)_i, \quad i = 1, \dots, n$$

$$\boldsymbol{\chi}_{i+n} = \hat{\mathbf{x}}_{t|t} - \left(\sqrt{(n + \lambda_{\text{ukf}})\mathbf{P}_{t|t}}\right)_i, \quad i = 1, \dots, n$$

where scaling parameter $$\lambda_{\text{ukf}} = \alpha_{\text{ukf}}^2(n + \kappa) - n$$, with spread parameter $$\alpha_{\text{ukf}} = 10^{-3}$$, weighting parameter $$\beta_{\text{ukf}} = 2.0$$ (optimal for Gaussian prior distributions), and parameter $$\kappa = 0$$.

The mean and covariance weighting coefficients are:

$$W_m^{(0)} = \frac{\lambda_{\text{ukf}}}{n + \lambda_{\text{ukf}}}, \quad W_c^{(0)} = W_m^{(0)} + (1 - \alpha_{\text{ukf}}^2 + \beta_{\text{ukf}})$$

$$W_m^{(i)} = W_c^{(i)} = \frac{1}{2(n + \lambda_{\text{ukf}})}, \quad i = 1, \dots, 2n$$

Each sigma point is propagated directly through the non-linear state transition function $$\mathbf{f}(\mathbf{x}, I_n)$$ without analytical linearization:

$$\boldsymbol{\chi}_{i, t+1|t} = \mathbf{f}(\boldsymbol{\chi}_{i, t|t}, I_{n,t})$$

$$\hat{\mathbf{x}}_{t+1|t} = \sum_{i=0}^{2n} W_m^{(i)} \boldsymbol{\chi}_{i, t+1|t}$$

$$\mathbf{P}_{t+1|t} = \sum_{i=0}^{2n} W_c^{(i)} \left[\boldsymbol{\chi}_{i, t+1|t} - \hat{\mathbf{x}}_{t+1|t}\right] \left[\boldsymbol{\chi}_{i, t+1|t} - \hat{\mathbf{x}}_{t+1|t}\right]^T + \mathbf{Q}$$

The predicted observation is computed by transforming each sigma point through the measurement model $$\mathcal{Y}_i = \mathbf{H} \boldsymbol{\chi}_{i, t+1|t} = -\chi_{i,1} - \chi_{i,2}$$:

$$\hat{y}_{t+1|t} = \sum_{i=0}^{2n} W_m^{(i)} \mathcal{Y}_i$$

$$P_{yy} = \sum_{i=0}^{2n} W_c^{(i)} (\mathcal{Y}_i - \hat{y}_{t+1|t})^2 + R$$

$$\mathbf{P}_{xy} = \sum_{i=0}^{2n} W_c^{(i)} \left[\boldsymbol{\chi}_{i, t+1|t} - \hat{\mathbf{x}}_{t+1|t}\right] (\mathcal{Y}_i - \hat{y}_{t+1|t})$$

The filter update follows:

$$\mathbf{K}_{t+1} = \mathbf{P}_{xy} P_{yy}^{-1}$$

$$\hat{\mathbf{x}}_{t+1|t+1} = \hat{\mathbf{x}}_{t+1|t} + \mathbf{K}_{t+1} (y_{t+1} - \hat{y}_{t+1|t})$$

$$\mathbf{P}_{t+1|t+1} = \mathbf{P}_{t+1|t} - \mathbf{K}_{t+1} P_{yy} \mathbf{K}_{t+1}^T$$

\begin{theorem}[Theorem 2.9 (High-Order Unscented State Reconstruction \& Curvature Invariance Theorem):] Under the deterministic sigma-point sampling scheme, the Unscented Kalman Filter captures the posterior mean $$\hat{\mathbf{x}}$$ and error covariance $$\mathbf{P}$$ accurately to third order ($$\mathcal{O}(\\|\Delta \mathbf{x}\\|^3)$$) of the Taylor series expansion for arbitrary prior probability densities. Furthermore, the UKF observer is immune to Jacobian singularity when physical capital investment collapses ($$I_{n,t} \to 0$$) or when ecological capacity approaches boundary exhaustion ($$R_{n,t} \to 0$$).

\textbf{Comparative Empirical Telemetry Under 2.0\% Gaussian Noise:}

\begin{table}[htbp]
\centering
\small
\caption{Comparative Observer Performance Under 2.0\% Gaussian Measurement Noise}
\label{tab:observer_comparison}
\begin{tabularx}{\textwidth}{l c c p{4.5cm}}
\toprule
\textbf{Estimation Metric} & \textbf{First-Order EKF (Thm 2.8)} & \textbf{Third-Order UKF (Thm 2.9)} & \textbf{Comparative Performance} \\
\midrule
Mean Relative Error ($R_n$) & 4.82\% & \textbf{4.78\%} & -0.04\% error reduction \\
Mean Squared Error ($R_a$) & 0.000625 & \textbf{0.000614} & -1.76\% MSE reduction \\
\textbf{Convergence Horizon} & 3 Quarters & \textbf{3 Quarters} & Equivalent fast lock \\
Trace Covariance $\text{Tr}(\mathbf{P}_{t|t})$ & 0.00038 & \textbf{0.00035} & -7.89\% tighter uncertainty \\
\textbf{Taylor Truncation Order} & $\mathcal{O}(\|\Delta \mathbf{x}\|^2)$ (1st order) & $\mathcal{O}(\|\Delta \mathbf{x}\|^4)$ (3rd order) & Exact 2nd-order curvature \\
Singularity Robustness ($I_n, R_n \to 0$) & Gradient divergence risk & \textbf{Completely robust} & Derivative-free stability \\
\bottomrule
\end{tabularx}
\end{table}

This establishes empirical and mathematical parity between the theoretical formulation and the verified implementation in work\_antiwork\_dynamics.py (UnscentedKalmanDecoupler) and test\_value\_physics.py (test\_unscented\_kalman\_decoupling\_observer). \end{theorem}

\section{Multi-Sector Empirical Calibration Suite}

\subsection{Empirical 56-Cell Value Matrices across 5 Essential Sectors}

We calibrate 56-cell ($7\text{ Planes} \times 8\text{ Stages}$) reference matrices across the 5 primary human provisioning sectors:

  
  

\begin{table}[htbp]
\centering
\small
\caption{Empirical Reference Standards and Distortion Quotients across 5 Essential Sectors}
\label{tab:matrix_2}
\begin{tabularx}{\textwidth}{l c c c p{3.8cm}}
\toprule
\textbf{Sector \& Reference Standard} & \textbf{Direct Diff ($WEST$)} & \textbf{Base Price} & \textbf{Primary $DQ$} & \textbf{Retail Toll $DQ$} \\
\midrule
1. Food (700g Bread Loaf $TS_0$) & 4.69 W (0.547 W dir) & \$5.00 AUD (\$1.15 base) & \textbf{0.26} (Stages 1--3) & \textbf{10.7} (Stage 8 Retail Toll) \\
2. Energy (1 kWh Microgrid) & 2.05 W (0.103 W dir) & \$0.285 AUD (\$0.075 base) & \textbf{0.24} (Stage 1 PV) & \textbf{15.1} (Stage 8 Gentailer Toll) \\
3. Water (1 kL Potable) & 2.80 W (0.140 W dir) & \$3.00 AUD (\$0.650 base) & \textbf{0.23} (Stage 1 Catch) & \textbf{13.7} (Stage 8 Utility Levy) \\
4. Shelter (1 Wk Housing) & 5.40 W direct & \$65.0 AUD/wk (\$21.0 base) & \textbf{0.31} (Stage 4 Build) & \textbf{18.3} (Stage 8 Mortgage Usury) \\
5. Care (1 Unit Health/Ed) & 3.20 W (1.75 W dir) & \$45.0 AUD (\$16.0 base) & \textbf{0.36} (Stage 4 Clinical) & \textbf{20.6} (Stage 8 Insurance Toll) \\
\bottomrule
\end{tabularx}
\end{table}

(Detailed concordance between 56-cell batch reference standards and single-unit continuous throughput coefficients is established in Appendix A).

\subsection{The Multi-Commodity Leontief-Value Tensor ($\mathbf{\Psi}_{\text{mesh}}$)}

Coupling all 5 sectors via the input-output technical coefficients matrix $\mathbf{A}$:  
$$\mathbf{A} = \begin{pmatrix} 0.000 & 0.001 & 0.002 & 0.500 & 0.100 \\ 0.050 & 0.000 & 0.450 & 15.000 & 1.200 \\ 0.002 & 0.001 & 0.000 & 0.800 & 0.050 \\ 0.001 & 0.0005 & 0.0005 & 0.000 & 0.020 \\ 0.002 & 0.001 & 0.001 & 0.050 & 0.000 \end{pmatrix}$$

  

Solving the Leontief inverse $(\mathbf{I} - \mathbf{A})^{-1}$ yields the total integrated embodied difficulty $\mathbf{d}{\text{integrated}} = \mathbf{d}{\text{direct}} \cdot (\mathbf{I} - \mathbf{A})^{-1}$:

  

  - \textbf{Food (1 Loaf):} Direct $0.547\text{ W} \longrightarrow$ Integrated $0.565\text{ W}$ (+3.2\% indirect exergy).
  - Energy (1 kWh): Direct $0.103\text{ W} \longrightarrow$ Integrated $0.109\text{ W}$ (+6.5\% indirect exergy).
  - Water (1 kL): Direct $0.140\text{ W} \longrightarrow$ Integrated $0.196\text{ W}$ (+40.1\% indirect exergy from electric pumping).
  - Shelter (1 Wk): Direct $5.40\text{ W} \longrightarrow$ Integrated $7.58\text{ W}$ (+40.4\% indirect exergy from materials/energy).
  - Care (1 Unit): Direct $1.75\text{ W} \longrightarrow$ Integrated $2.10\text{ W}$ (+19.9\% indirect exergy from facilities).

\section{Stock-Flow Consistent (SFC) Dynamic Micro-Economy Simulation}

\subsection{Empirical Macroeconomic Calibration}

We construct a 156-week dynamic Stock-Flow Consistent (SFC) model of a representative Australian regional shire (600 citizens, 200 households, 10 broadacre cropping enterprises managing 4,000 ha) calibrated strictly against official microdata from ABARES, ABS, RBA, ACCC, and AEMO.

  

Following Godley \& Lavoie (2007), all transaction flows satisfy exact zero-sum balance sheet closure ($\sum \text{Rows} \equiv 0, \sum \text{Columns} \equiv 0$), bifurcating current account interest flows ($i \cdot D$) from capital account linear amortization ($P_{\text{linear}}$).

\subsection{Monetary Policy Working Capital Cost Channel}

Under compound macroeconomic shocks (RBA cash rate hike from 0.10\% to 4.35\%, diesel fuel increase from \$1.45 to \$2.25/L, 30-week 40\% drought yield shock):

  

$$\text{Unit Production Cost: } c_t = \frac{w L_t + p_e E_t + i_t D_{\text{working}}}{Y_t}$$

  

Grounded in the working capital cost channel (Barth \& Ramey 2001; Ravenna \& Walsh 2006; Tillmann 2008), rate hikes directly increase upstream interest expenses from \$10,200 to \$24,900/week across the 10 broadacre farms. Farmers cannot pass these costs downstream due to monopsonistic duopoly wholesale gatekeeping ($DQ_{\text{farmgate}} = 0.26$). In contrast, retail duopolies defend their \~28\% gross margins against inelastic consumer demand ($\epsilon \to 0$), driving retail bread prices from \$4.60 to \$5.65 AUD.

\section{Non-Usurious Credit Clearing: Theorem 5 \& The Tri-Fold Ledger}

\subsection{Theorem 5 (Statutory Debt Bifurcation)}

\textbf{Theorem 5:} To prevent exponential debt compounding from inducing physical substrate collapse, credit must be bifurcated into Class A (Productive Capital) and Class B (Discretionary Credit). For Class A facilities:  
$$r_{\text{compound}} \equiv 0.0\%, \quad \text{Amortization}(t) = \frac{C_{\text{principal}}}{T_{\text{term}}}$$

\subsection{The Tri-Fold Cryptographic Ledger \& PoD Consensus}

  - \textbf{Book 3 (Reality Ledger):} Cryptographically logs physical exergy ($P_e$) and labor difficulty ($P_b$) into deterministic Proof-of-Difficulty block headers (struct PoD\_BlockHeader with bytes65 guild\_multisig[3]).
  - \textbf{Book 2 (Public Sovereign Ledger):} Issues non-inflationary sovereign credits pegged 1:1 to verified $WEST$ tokens to fund strategic physical buffers ($TS_0$ Silo, Microgrid, Water Dam).
  - \textbf{Book 1 (Private Mutual Credit Ledger):} Clears bilateral member trade at 0.0\% compound interest.

\subsection{1,000-Run Monte Carlo Stress Testing}

Across 1,000 independent 156-week stochastic shock iterations:

  

  - \textbf{Neoclassical Baseline:} Mean farm debt expanded from \$18.0M to \$19.2M AUD (VaR 95\% = \$20.3M), FMD reserves completely drained, and cumulative reality debt reached \$4.93M AUD in unhedged soil/capital wear.
  - Value Physics Protocol: Amortized \$2.16M of principal debt linearly, accumulated \$6.18M AUD in regenerative capital, maintained 100\% solvency ($P_{\text{default}} = 0.000$), and stabilized local bread at \$1.15 AUD ($DQ = 1.000, \\Delta R = \$0.00$).

\section{Conclusions \& Policy Implications}

This research demonstrates that economic value is fundamentally anchored in the thermodynamic and biological reproduction cost of life. The Universal Price Equation unifies physical exergy, metabolic difficulty, and institutional toll extraction into a rigorous, predictive tensor framework.

  

By enacting Statutory Debt Bifurcation (Theorem 5) and establishing Regional Sovereign Mutual Credit institutions, economies can permanently decouple essential human provisioning from financialized interest compounding, eliminating artificial cost-push inflation and securing long-term ecological and societal resilience.

\section{Author Epistemic Statement, CRediT, Data \& Code Availability}

\subsection*{Author Epistemic Statement \& AI Disclosure}
The conceptual ontology, foundational equations, and analytical architecture of Value Physics represent the original intellectual work of Jarrod Lyndsay Hamilton. The author holds no formal academic credentials in economics; theoretical connections were established through self-directed study. Due to memory impairments, the author may not actively retain specific textual details or citations from the research process; readers are encouraged to evaluate the mathematical mechanics, computational solvers, and empirical calibrations directly on their own intrinsic merits. AI tooling (Google Gemini / Gemini Spark) was utilized under the author's direct guidance for computational implementation, code verification, and \LaTeX\ formatting.

\subsection{CRediT Authorship Statement}
\textbf{Jarrod Lyndsay Hamilton:} Conceptualization, Methodology, Formal Proofs, Investigation, Data Curation, Writing -- Original Draft, Visualization.

\subsection{Data Availability Statement}

All baseline empirical datasets used in the SFC dynamic simulation and 56-cell matrix calibrations are derived from publicly accessible official statistical repositories:

1.  1\. Australian Bureau of Agricultural and Resource Economics and Sciences (ABARES) Agricultural Commodity Microdata.
2.  2\. Australian Bureau of Statistics (ABS) Consumer Price Index and Input-Output Tables (Cat. 5209.0.55.001). 
3.  3\. Reserve Bank of Australia (RBA) Historical Cash Rate Target Series (Table F1.1).
4.  4\. Australian Competition and Consumer Commission (ACCC) Supermarket Inquiry Reports (2024–2026).
5.  5\. Australian Energy Market Operator (AEMO) Regional NSW Electricity Generation Telemetry.

\subsection{Code Availability Statement}

The complete replication codebase—including value\_physics\_engine.py, multi\_village\_simulation.py, finite\_boundary\_stress\_test.py, and the automated test suite test\_value\_physics.py—is available in the open-source replication repository with pinned random seeds for full numerical reproducibility.

\subsection{Conflict of Interest Declaration}

The authors declare that they have no known competing financial interests or personal relationships that could have appeared to influence the work reported in this paper.





\section*{Appendix A: Physical Input-Output Table (PIOT) \& Batch-to-Unit Concordance}
\addcontentsline{toc}{section}{Appendix A: Physical Input-Output Table (PIOT) \& Batch-to-Unit Concordance}

To bridge the 56-cell empirical engineering matrix standards (which represent batch-level production runs for observational accounting) with the continuous single-unit technical coefficients matrix $\mathbf{A}$, Table A.1 defines the explicit scale concordance factors:

\subsubsection*{Table A.1: Batch-to-Unit Multi-Sector Scale Concordance}

\begin{table}[htbp]
\centering
\small
\caption{Batch-to-Unit Multi-Sector Scale Concordance}
\label{tab:scale_concordance}
\begin{tabularx}{\textwidth}{l p{4.2cm} p{4.2cm} c r r}
\toprule
\textbf{Sector} & \textbf{Empirical Batch (56-Cell)} & \textbf{Unit Standard (Leontief $\mathbf{A}$)} & \textbf{Scale $k_i$} & \textbf{Direct Diff} & \textbf{Base Cost ($P_H$)} \\
\midrule
1. Food & 1 Artisan Batch (8.57 Loaves) & 1 Industrial Loaf (700g) & $8.57\times$ & 0.547 WEST & \$5.00 $\to$ \$1.15 AUD \\
2. Energy & 1 Household Daily ($20.0\text{ kWh}$) & 1 Kilowatt-Hour ($1.0\text{ kWh}$) & $20.0\times$ & 0.103 WEST & \$5.70 $\to$ \$0.075 AUD \\
3. Water & 1 Reticulation Block ($20.0\text{ kL}$) & 1 Kilolitre ($1.0\text{ kL}$) & $20.0\times$ & 0.140 WEST & \$60.00 $\to$ \$0.650 AUD \\
4. Shelter & 1 Resident Weekly Allocation & 1 Resident-Week Housing & $1.00\times$ & 5.400 WEST & \$65.00 $\to$ \$21.00 AUD \\
5. Care & 1 Team Session ($1.83\text{ hrs}$) & 1 Patient Hour ($1.0\text{ hr}$) & $1.83\times$ & 1.750 WEST & \$82.50 $\to$ \$16.00 AUD \\
\bottomrule
\end{tabularx}
\end{table}

For any sector $i$, continuous direct difficulty is computed as $d_{\text{direct}, i} = \frac{D_{\text{batch}, i}}{k_i}$, guaranteeing mathematical continuity across micro-accounting and macroeconomic input-output inverses.



\section*{Appendix B: Microeconomic Coercion Kernel Derivation}
\addcontentsline{toc}{section}{Appendix B: Microeconomic Coercion Kernel Derivation}

We derive the demand price elasticity collapse $\epsilon(v_{\text{rel}}) \to 0$ and rentier margin capture $R_a \to 1.0$ from foundational consumer and bargaining theory.

\subsection{Stone-Geary Subsistence Exhaustion}

Consider an agent maximizing a non-homothetic Stone-Geary utility function (Geary 1950; Stone 1954) over a subsistence necessity good $x$ (with non-negotiable biological floor $\gamma_x > 0$) and a composite discretionary good $y$:  
$$\max_{x, y} U(x, y) = (x - \gamma_x)^{\alpha} (y - \gamma_y)^{1-\alpha} \quad \text{subject to } P_x x + P_y y = M$$  
Where $M$ is total nominal expenditure capacity. Solving the Lagrangian yields the Marshallian demand function for subsistence good $x$:  
$$x(P_x, P_y, M) = \gamma_x + \frac{\alpha (M - P_x \gamma_x - P_y \gamma_y)}{P_x}$$  
The uncompensated (Marshallian) price elasticity of demand $\epsilon_x = -\frac{\partial \ln x}{\partial \ln P_x}$ is analytically derived as:  
$$\epsilon_x = -\frac{\partial x}{\partial P_x} \frac{P_x}{x} = \frac{(1 - \alpha) P_x \gamma_x + \alpha P_y \gamma_y}{P_x \gamma_x + \alpha (M - P_x \gamma_x - P_y \gamma_y)}$$  
Under existential survival urgency ($U \to \infty$), discretionary consumption collapses ($\gamma_y \to 0, y \to 0$) and total expenditure approaches the subsistence boundary ($M \to P_x \gamma_x$). In this limit:  
$$\lim_{M \to P_x \gamma_x} \epsilon_x = \frac{(1 - \alpha) P_x \gamma_x}{P_x \gamma_x} = 1 - \alpha$$  
For biological survival staples (where subsistence priority $\alpha \to 1.0$), elasticity collapses strictly to zero:  
$$\lim_{\alpha \to 1.0} \epsilon_x = 0$$

\subsection{Rubinstein Alternating-Offers Urgency Asymmetry}

Consider an alternating-offers bargaining game (Rubinstein 1982) between Consumer A (facing existential urgency $U_A$) and Supplier B (holding store inventory with zero immediate urgency $U_B = 0$).  
Let the discount factors per time interval $\Delta t$ be:  
$$\delta_A = e^{-r_A \Delta t}, \quad \delta_B = e^{-r_B \Delta t}$$  
Where subjective discount rate $r_A = r_0 (1 + U_A)$ reflects the biological cost of waiting (e.g. starvation, dehydration). As urgency $U_A \to \infty$, $r_A \to \infty \implies \delta_A \to 0$.  
The unique Subgame Perfect Equilibrium (SPE) share of surplus captured by Supplier B is given by:  
$$s_B^* = \frac{1 - \delta_A}{1 - \delta_A \delta_B}$$  
Taking the asymptotic limit as consumer impatience approaches infinity ($\delta_A \to 0$):  
$$\lim_{\delta_A \to 0} s_B^* = \frac{1 - 0}{1 - 0} = 1.000$$  
Supplier B captures 100\% of the economic surplus. The resulting equilibrium transaction price is:  
$$P^* = P_H + s_B^* (P_{\text{max}} - P_H) \xrightarrow{s_B^* \to 1.0} P_{\text{max}} = \frac{P_H}{1 - R_a}$$  
Where rentier enclosure fraction $R_a \to 1.0$ and Lorentz coercion factor $\gamma(v_{\text{rel}}) \to \infty$ naturally emerge at the survival boundary.


\bibliographystyle{elsarticle-harv}
\bibliography{references_vpe21}

\end{document}
